\begin{textsample}{NEG Dim 5 – pi – Score 1.00 – pi/pi\_ef\_1.txt}  \label{ex:f5_neg_058}
Ensino Fundamental LÍNGUA PORTUGUESA MARCO LEGAL Ao elaborar a Base Nacional Comum Curricular, o Brasil promove uma forma de unificação do ensino no País.
Era uma necessidade que tínhamos como educadores : não sentir diferenças na forma de ensinar.
Para que esse intuito seja atendido, é imprescindível que todos os estados se equiparem com as discussões e metas da Base Nacional Comum Curricular.
Conforme definido pela Lei de Diretrizes e Bases da Educação Nacional ( LDB, Lei nº 9.394/1996 ), a BNCC deve nortear os currículos dos sistemas e redes de ensino das Unidades Federativas, como também as propostas pedagógicas de todas as escolas públicas e privadas de Educação Infantil, Ensino Fundamental e Ensino Médio, em todo o Brasil.
Nesse cenário, apresentamos o Currículo Estadual de Língua Portuguesa, elaborado em grupos de discussão por professores da educação básica, gestores e Dirigentes Municipais de Educação numa perspectiva de educação integral para os estudantes piauienses.
Alinhamo -nos à BASE e, dessa forma, consideramos a língua como um sistema capaz de produzir relações e é nesse sentido que ela assume função de inclusão e/ou exclusão.

% matched lemmas: 
\end{textsample}
